\documentclass[11pt]{article}
\usepackage[margin=1in]{geometry}
\usepackage{amsmath, amssymb}
\usepackage{xcolor}
\usepackage{enumitem}
\usepackage{titlesec}
\usepackage[colorlinks=true, linkcolor=blue!50!black, urlcolor=blue!50!black, citecolor=blue!50!black]{hyperref}
\usepackage{fancyhdr}

\definecolor{codeblue}{HTML}{213A87}
\definecolor{accentred}{HTML}{BF1F1F}

\titleformat{\section}{\Large\bfseries\color{codeblue}}{\thesection}{0.8em}{}

\pagestyle{fancy}
\fancyhf{}
\lhead{\small Precomps Bibliography --- PAT-scan Literature}
\rhead{\small \thepage}
\renewcommand{\headrulewidth}{0.4pt}

\setlength{\parindent}{0pt}
\setlength{\parskip}{4pt}

\title{Precomps Bibliography \\[4pt]
  \large User-Curated Pre-Computation Set (Stack 3)}
\author{Claude Opus 4.7}
\date{May 31, 2026}

\newcommand{\bibentry}[6]{%
  \noindent\textbf{[#1]}~#2.~\textit{#3}.~#4, #5.~%
  \href{https://doi.org/#6}{\texttt{doi:#6}}\par\vspace{6pt}
}

\begin{document}
\maketitle
\thispagestyle{fancy}

\begin{center}
\small\itshape
Four papers enriched via OpenAlex from the flat \texttt{precomps/} folder. \\
Three overlap with the shortlisted stack; one (JAX-SSO 2024) is unique to this folder.
\end{center}

\vspace{1em}
\hrule
\vspace{1em}

%% ============================================================
\section*{Summary}

\begin{tabular}{@{}lrl@{}}
\textbf{Item} & \textbf{Value} & \\
\hline
Total papers          & 4 & \\
Year range            & 2003--2024 & \\
Shared with shortlisted & 3 & (Mei 2017, Konofagou 2004, Feng 2022) \\
Unique to precomps    & 1 & (Wu 2024 --- JAX-SSO) \\
Open-access           & 3 & (gold $\times$1, green $\times$2) \\
\hline
\end{tabular}

\vspace{1em}

%% ============================================================
\section{Precomputation References}

\bibentry{1}{Konofagou, E.\,E.\ and Harrigan, T.\,P.}{Palpation tomography --- a new technique for modulus estimation in elastography}{IEEE Ultrasonics Symposium Proceedings}{2003/2004}{10.1109/ultsym.2003.1293486}

\textit{Columbia University; Exponent (United States).}\quad
Multiple smaller quasi-static loads (inspired by clinical palpation) replace a single whole-boundary load, increasing the measurement-to-parameter ratio and reducing noise sensitivity. \emph{Also appears in shortlisted/Statics.}

\bibentry{2}{Mei, Y., Wang, S., Shen, X., Rabke, S., and Goenezen, S.}{Mechanics based tomography: a preliminary feasibility study}{Sensors (MDPI)}{2017}{10.3390/s17051075}

\textit{Texas A\&M University.}\quad Cited by 19. \quad Open access (gold). \quad
Coins the term \emph{Mechanics Based Tomography} (MBT): force sensors plus DIC-measured boundary displacements --- no interior measurements --- reconstruct shear-modulus maps of stiff inclusions in a soft background. \emph{Also appears in shortlisted/Statics.}

\bibentry{3}{Feng, B.\,T., Ogren, A.\,C., Daraio, C., and Bouman, K.\,L.}{Visual vibration tomography: estimating interior material properties from monocular video}{IEEE/CVF Conference on Computer Vision and Pattern Recognition (CVPR)}{2022}{10.1109/cvpr52688.2022.01575}

\textit{California Institute of Technology.}\quad Cited by 12.\quad Open access (green). \quad
Estimates spatially-varying Young's modulus and density throughout a 3D object from a single monocular video, by extracting image-space modes from sub-pixel surface vibration and inverting to material properties. \emph{Also appears in shortlisted/Vibrations.}

\bibentry{4}{Wu, G.}{JAX-SSO: differentiable finite element analysis solver for structural optimization and seamless integration with neural networks}{arXiv (Cornell University)}{2024}{10.48550/arxiv.2407.20026}

\textit{Single-author preprint.}\quad Open access (green). \quad
Differentiable FEA solver built on Google's JAX. Uses the adjoint method plus automatic differentiation for sensitivity calculation in structural optimization (shape, size, topology) and integrates natively with neural networks for physics-informed training. GPU-accelerated. \emph{Unique to the precomps folder --- this is the only paper not in the shortlisted bibliography.}

\vspace{2em}
\hrule
\vspace{0.5em}

{\footnotesize\itshape
Generated by the Stack 3 cap-stone skill (\texttt{/identify-papers-user}) over the flat folder at \\
\texttt{paper-stack-three-test-two/precomps/}. Source: OpenAlex (resolution mode: search) on May 31, 2026.
}

\end{document}
